% !PRINTED 491
\weberpage{491}

Wenn man $u = v$ setzt, so wird

\[
A = 2, \quad B = u^2 + \frac{2}{u^2},
\]
und wenn $u = f(\sqrt{-m})$, $m = 7, 19, 23$ gesetzt wird, so erhält man aus \S\,128, (6), (8), \S\,131, (1) durch Elimination von $u$ die Gleichungen für $B$:

\[
\begin{aligned}
m &= 7, & B - 3 &= 0, \\
m &= 19, & B^2 - 6B^2 + 10B - 6 &= 0, \\ % [sic] printed first term is $B^2$; the sentence below and (7) require $B^3$
m &= 23, & B^3 - 5B^2 + 4B - 1 &= 0.
\end{aligned}
\]

Die beiden letzten kubischen Gleichungen sind, da sie keine rationale Wurzel haben, irreducibel, und daraus ergibt sich ohne Rechnung die gesuchte Modulargleichung (3):

\weq{7}{A - 2 = (B-3)^2\,(B^3 - 6B^2 + 10B - 6)^2\,(B^3 - 5B^2 + 4B - 1).}

In der Abhandlung: \glqq Ein Beitrag zur Transformationstheorie der elliptischen Funktionen mit einer Anwendung auf Zahlentheorie\grqq{}, Math. Annalen 43, 185, habe ich auf demselben Wege auch noch die Transformationsgleichung für den 47sten Transformationsgrad abgeleitet.

\begin{center}
\S\,133. \gesperrt{Die Resolventen 7ten und 11ten Grades für den 7ten und 11ten Transformationsgrad.}
\end{center}

Wir haben in \S\,82 gesehen, daß für den 7ten und 11ten Transformationsgrad Resolventen der Grade 7 und 11 existieren. Auch bei der Bildung dieser Gleichungen kann die Theorie der komplexen Multiplikation nützliche Dienste leisten.

Wir betrachten, wenn $n = 7$ oder $= 11$ ist, die Transformationsgleichung

\[
F_n(u, v) = 0,
\]
deren Wurzeln, wenn $u = f(\omega)$ gesetzt wird,

\weq{1}{v_\infty = f(n\,\omega), \quad v_c = f\left(\frac{\omega+c}{n}\right)}

sind, wobei $c$ als Index von $v$ nach dem Modul $n$ genommen werden kann, während es unter dem Zeichen $f$ durch 48 teilbar vorausgesetzt werden muß.

Diese Gleichungen sind nach \S\,73:

\weq{2}{n = 7, \quad v^8 - u^7v^7 + 7u^4v^4 - 8\,uv + u^8 = 0,}

\weq{3}{n = 11, \quad v^{12} - u^{11}v^{11} + 11\,u^9v^9 - 44\,u^7v^7 + 88\,u^5v^5 - 88\,u^3v^3 + 32\,uv + u^{12} = 0.}
